
{\sffamily\bfseries\Large\color{accent}Routes through the guide\par}\vspace{12pt}
Plates VI.1 through VI.42 identify worked pages independently of PDF pagination. The references on Plates VI.1--VI.28 use the theorem numbering of Traditional Part I. References from Plate VI.29 onward name the corresponding strengthened result in Traditional Part II. These references distinguish preserved manuscript branches rather than merging their result counts.
\begin{center}\small
\begin{tabularx}{\linewidth}{@{}p{1.6in}Yl@{}}\toprule
Route & Mathematical question & Plates\\\midrule
Complete transfer & Which coordinates and modes reconstruct the same evolution? & 1--6, 29\\
Probability and evidence & What survives averaging, aggregation, or an observation? & 7--14, 26, 32, 34--36\\
Spectral evolution & How do rates, gaps, fractional clocks, and Gaussian modes interact? & 15--17, 25, 33, 41\\
Green and path formulas & How do grid units and weighting conventions enter? & 18--20\\
Eliminating hidden states & What memory returns, and how can it be compressed? & 21--24, 37--40\\
Passing to limits & What additional conditions give convergence and closure? & 27, 30--31\\
Evidence & Which calculations and formal records support which claims? & 28, 42\\\bottomrule
\end{tabularx}
\end{center}
\head{Three distinctions to keep visible}
A coordinate decomposition retains the whole vector if all components are kept. An aggregation map usually loses information. Elimination recovers exact retained dynamics only after adding a memory term and any hidden-state initial forcing.

An exact example checks its stated finite model. A theorem provides an argument under quantified assumptions. A Lean record certifies its declared formal statement under the recorded environment. None of these automatically certifies the other two.
\head{Preservation and revision}
The first twenty-eight worked pages develop the finite models; fourteen further pages connect them to the extended results. Forty editable vector figures illustrate these constructions. Their source files and retained generation records accompany the mathematical text and its formal correspondence.

The inherited formal baseline records 33 scoped results, 341 supporting theorem declarations, and 47 proof modules. The strengthened traditional branch is a separate written-proof collection. Its broader content is not added to the formal denominator.

Traditional Part III retains complementary results with an R prefix: reversibility and Dirichlet forms (R2), prior disintegrations and filtering (R3), potential perturbation and coarse Feynman--Kac transfer (R4), and finite realizations (R5). Related models appear on Plates VI.5--VI.8, Plate VI.11, Plate VI.13, Plate VI.20, Plate VI.26, and Plates VI.39--VI.40. The bounded-operator reciprocal result migrated from Chapter IV is R7.2. Plate VI.25 and Plate VI.41 explain related finite matrix calculus. These models do not replace the supplementary hypotheses.
\takeaway{Choose a question, work through its small model, then check the hypotheses in the corresponding traditional proof and formal contract.}
